\documentclass[aspectratio=169]{beamer}
\usetheme{Madrid}
\usecolortheme{beaver}
\usepackage[utf8]{inputenc}
\usepackage[T1]{fontenc}
\usepackage{tikz}
\usepackage{amssymb}
\usepackage{graphicx}
\setbeamertemplate{navigation symbols}{}
\setbeamercolor{title}{fg=white,bg=blue!55!black}
\setbeamercolor{frametitle}{fg=white,bg=blue!45!black}
\setbeamercolor{block title}{bg=blue!15,fg=black}
\setbeamercolor{block body}{bg=blue!5,fg=black}

\title{M1 — From Blocks to Text}
\subtitle{print, input, variables, types — the G7 skills, now in real Python}
\author{ByteForge --- Grade 8 --- App Builder}
\date{}

\begin{document}

\begin{frame}
    \titlepage
\end{frame}

\begin{frame}{Today's Plan}
    \begin{itemize}
        \item Say Hello
        \item About Me
        \item Type Detective
        \item Ask the User
        \item Temperature Converter
        \item[$\bigstar$] BOSS --- M1 Capstone
    \end{itemize}
\end{frame}

\begin{frame}{print() — your first real line of code}
    \Large
    \texttt{print()} displays whatever you put inside its brackets. This exact line is traditionally the very first thing anyone writes in a new programming language.
\end{frame}


\begin{frame}{Variables and f-strings}
    \Large
    A variable is still a named box holding a value — just written as \texttt{name = "Alex"} instead of snapped together. \texttt{f"My name is \{name\}"} drops a variable directly INSIDE a string.
\end{frame}


\begin{frame}{Every value has a TYPE}
    \Large
    \texttt{int} (7), \texttt{float} (7.5), \texttt{str} ("hi"), \texttt{bool} (True/False) — Python tracks the type of every value automatically. \texttt{type(x)} reveals it.
\end{frame}


\begin{frame}{input() always returns TEXT}
    \Large
    \texttt{input()} ALWAYS hands back text — typing "25" gives you the STRING "25", not the number 25. \texttt{int(age\_text)} converts it into a real number.
\end{frame}


\begin{frame}{A real formula, not just a print statement}
    \Large
    \texttt{celsius * 9 / 5 + 32} — Python does x and / before + and -, exactly like normal maths rules. A genuine, widely-used real-world formula.
\end{frame}

\begin{frame}{Real-World Example}
    \Large
    Every weather app converting between Celsius and Fahrenheit runs this EXACT formula. Every sign-up form that asks your age and immediately validates it depends on the same input()-returns-text, then-convert pattern.
\end{frame}


\begin{frame}{Image}
    \begin{center}
    \includegraphics[width=0.5\textwidth]{images/python.png}

    \large The real language you're about to start typing.

    \small Python logo, CC BY 2.5, via Wikimedia Commons
    \end{center}
\end{frame}

\begin{frame}{Let Them Discover}
    \Large
    Type \texttt{type(input("Enter a number: "))} for yourself conceptually — before running anything, guess: will Python say the result is a number, or text? Most students guess wrong the first time.
\end{frame}

\begin{frame}{Your Turn}
    \Large
    Open the student map, find \textbf{M1 — From Blocks to Text}, and work through it
    hands-on --- everything you just saw, they now do themselves.
\end{frame}

\end{document}
